\subsection{ASSOCIATED SUBGROUPS}

Let \(\mathbf{G}\) be an arbitrary subgroup (closed or not) of \(\mathbf{R}^n\) . Consider the set \(\mathbf{G}^*\) of points \(\boldsymbol{u} = (u_i) \in \mathbf{R}^n\) such that, for all \(\boldsymbol{x} = (x_i) \in \mathbf{G}\) , the number \((\boldsymbol{u}| \boldsymbol{x}) = \sum_{i=1}^{n} u_i x_i\) is an integer. It is immediately seen that \(\mathbf{G}^*\) is a subgroup of \(\mathbf{R}^n\) ; it is called the subgroup associated with \(\mathbf{G}(*)\) . If \(\mathbf{G}\) and \(\mathbf{H}\) are two subgroups of \(\mathbf{R}^n\) such that \(\mathbf{H} \subset \mathbf{G}\) , it is clear that \(\mathbf{G}^* \subset \mathbf{H}^*\) .

PROPOSITION 4. The subgroup \(\mathbf{G}^*\) associated with a subgroup \(\mathbf{G}\) of \(\mathbf{R}^n\) is closed, and we have \((\overline{\mathbf{G}})^* = \mathbf{G}^*\) .

For each \(x \in G\) , let \(f_x(u)\) denote \((u|x)\) ; \(f_x\) is a linear form, therefore continuous. Since \(G^*\) is the intersection of the sets \(f_x^{-1}(Z)\) as \(x\) runs through \(G\) , and since each of these sets is closed, it follows that \(G^*\) is closed. On the other hand, if \(u \in G^*\) , we have \((u|x) \in Z\) for all \(x \in G\) , and therefore, as \(Z\) is closed in \(R\) , \((u|y) \in Z\) for all \(y \in \overline{G}\) ; thus \(u \in (G)^*\) . But we have \((\overline{G})^* \subset G^*\) (since \(G \subset \overline{G}\) ), so that \((\overline{G})^* = G^*\) .

Consider the structure of \(\mathbf{G}^*\) when \(\mathbf{G}\) is closed. By Corollary 1 to Theorem 2 of no. 2, there exists a base \((\pmb{a}_i)_{\mathbf{1} \leqslant i \leqslant n}\) of \(\mathbf{R}^n\) such that \(\mathbf{G}\) coincides with the set of points

\[
\boldsymbol {x} = \sum_ {i = 1} ^ {p} t _ {i} \boldsymbol {a} _ {i} + \sum_ {j = p + 1} ^ {p + q} n _ {j} \boldsymbol {a} _ {j},
\]

where the \(t_i\) take all real values and the \(n_j\) take all integer values. Hence \((\boldsymbol{u}|\boldsymbol{x})\) is an integer for all these points \(\boldsymbol{x}\) if and only if \((\boldsymbol{u}|\boldsymbol{a}_i) = 0\) for \(1 \leqslant i \leqslant p\) and \((\boldsymbol{u}|\boldsymbol{a}_i)\) is an integer for \(p + 1 \leqslant i \leqslant p + q\) . Let us denote by \((\boldsymbol{a}_i')_{1 \leqslant i \leqslant n}\) the basis of \(\mathbf{R}^n\) such that \((\boldsymbol{a}_i'|\boldsymbol{a}_j) = 0\) for \(i \neq j\) and

\[
\left(\boldsymbol {a} _ {i} ^ {\prime} \mid \boldsymbol {a} _ {i}\right) = \mathrm{I} \quad \text { for } \quad \mathrm{I} \leqslant i \leqslant n
\]

{[}the basis `` dual '' to \((a_{i})\) {]}; if we put \(u = \sum_{i=1}^{n} u_{i} a_{i}^{\prime}\) , it is clear that the points \(u \in G^{*}\) are characterized by the following conditions: \(u_{i} = 0\) for \(i \leqslant i \leqslant p\) , and \(u_{i} \in Z\) for

\[
p + \mathrm{i} \leqslant i \leqslant p + q;
\]

hence \(G^{*}\) is the direct sum of a vector subspace W with a basis consisting of the \(a_{i}^{\prime}\) such that \(p + q + i \leqslant i \leqslant n\) , and a discrete subgroup generated by the \(a_{i}^{\prime}\) such that \(p + i \leqslant i \leqslant p + q\) . In other words:

PROPOSITION 5. For every closed subgroup G of \(\mathbf{R}^n\) ,

\[
r (\mathbf {G} ^ {*}) = n - d (\mathbf {G}) \quad a n d \quad d (\mathbf {G} ^ {*}) = n - r (\mathbf {G}).
\]

Let us apply the same reasoning to \(\mathbf{G}^*\) . Observing that the basis dual to \((\pmb{a}_i')\) is \((\pmb{a}_i)\) , we see that:

PROPOSITION 6. For every subgroup \(\mathbf{G}\) of \(\mathbf{R}^n\) , we have \((\mathbf{G}^*)^* = \overline{\mathbf{G}}\) .

COROLLARY. A point x lies in the closure of a subgroup G of \(R^{n}\) if and only if \((u|x)\) is an integer for all \(u \in R^{n}\) such that \((u|y)\) is an integer for all \(y \in G\) .

Let us apply this characterization of points lying in the closure of a subgroup \(\mathbf{G}\) to the case of the subgroup \(\mathbf{G}\) generated by the \(n\) vectors \(\boldsymbol{e}_j\) of the canonical basis ( \(\mathrm{i} \leqslant j \leqslant n\) ) and by an arbitrary number \(m\) of points \(\boldsymbol{a}_i\) ( \(\mathrm{i} \leqslant i \leqslant m\) ) of \(\mathbf{R}^n\) . To say that \((\boldsymbol{u}|\boldsymbol{e}_j)\) is an integer for \(\mathrm{i} \leqslant j \leqslant n\) means that the \(n\) coordinates of \(\boldsymbol{u}\) are integers; therefore:

PROPOSITION 7 (Kronecker). Let \(\boldsymbol{a}_i = (a_{ji})\) ( \(i \leqslant i \leqslant m\) , \(i \leqslant j \leqslant n\) ) be \(m\) points of \(\mathbf{R}^n\) and let \(\boldsymbol{b} = (b_j)\) ( \(i \leqslant j \leqslant n\) ) be a point of \(\mathbf{R}^n\) . In order that, for each \(\varepsilon > 0\) , there exist \(m\) integers \(g_i\) ( \(i \leqslant i \leqslant m\) ) and \(n\) integers \(p_j\) ( \(i \leqslant j \leqslant n\) ) such that

\[
\left| q _ {1} a _ {1 j} + q _ {2} a _ {2 j} + \dots + q _ {m} a _ {m j} - p _ {j} - b _ {j} \right| \leqslant \varepsilon \quad (\mathrm{I} \leqslant j \leqslant n)
\]

it is necessary and sufficient that, for each finite sequence \((r_j)\) ( \(1 \leqslant j \leqslant n\) ) of \(n\) integers such that the \(m\) numbers \(\sum_{j=1}^{n} a_{ij} r_j\) ( \(1 \leqslant i \leqslant m\) ) are all integers, the number \(\sum_{j=1}^{n} b_j r_j\) should also be an integer.

COROLLARY 1. In order that, for each \(\boldsymbol{x} = (x_{j}) (\mathrm{I} \leqslant j \leqslant n)\) and each \(\varepsilon > 0\) , there exist m integers \(q_{i} (\mathrm{I} \leqslant i \leqslant m)\) and n integers

\(p_j\) (1 \(\leqslant j\leqslant n\) ) such that

\[
\left| q _ {1} a _ {1 j} + q _ {2} a _ {2 j} + \dots + q _ {m} a _ {m j} - p _ {j} - x _ {j} \right| \leqslant \varepsilon \quad (\mathrm{I} \leqslant j \leqslant n)
\]

it is necessary and sufficient that there exist no finite sequence \((r_j)\) of \(n\) integers, not all zero, such that each of the \(m\) numbers \(\sum_{j=1}^{n} a_{ij} r_j\) is an integer. For if \(G\) is dense in \(\mathbf{R}^n\) , that is if \(\overline{\mathbf{G}} = \mathbf{R}^n\) , then \(\mathbf{G}^* = \{o\}\) , and conversely.

In particular \((m = \mathrm{i})\) :

COROLLARY 2. Let \(\theta_{1},\theta_{2},\ldots ,\theta_{n}\) be \(n\) real numbers. In order that, given any \(n\) real numbers \(x_{1},x_{2},\ldots ,x_{n}\) and a real number \(\varepsilon >0\) , there should exist an integer \(q\) and \(n\) integers \(p_j\) such that

\[
\left| q \theta_ {j} - p _ {j} - x _ {j} \right| \leqslant \varepsilon \quad (\mathrm{I} \leqslant j \leqslant n)
\]

it is necessary and sufficient that there exist no relation of the form \(\sum_{j=1}^{n} r_j \theta_j = h\) , where the \(r_j\) are \(n\) integers not all zero, and \(h\) is an integer {[}which implies, in particular, that the \(\theta_j\) and the ratios \(\theta_j / \theta_k\) ( \(j \neq k\) ) must be irrational{]}.

We may interpret this result as follows: for each integer \(q \in Z\) , let \(x_{q}\) denote the point with coordinates \(q\theta_{j} - [q\theta_{j}]\) ( \(i \leqslant j \leqslant n\) ); then Corollary 2 gives a necessary and sufficient condition for the set of points \(x_{q}\) to be dense in the cube which is the product of n intervals {[}0, i{]} in the factors of \(R^{n}\) .

PROPOSITION 8. If \(G_{1}\) , \(G_{2}\) are any closed subgroups of \(R^{n}\) , we have

\[
(\mathrm{G} _ {1} + \mathrm{G} _ {2}) ^ {*} = \mathrm{G} _ {1} ^ {*} \cap \mathrm{G} _ {2} ^ {*} \quad a n d \quad (\mathrm{G} _ {1} \cap \mathrm{G} _ {2}) ^ {*} = \mathrm{G} _ {1} ^ {*} + \mathrm{G} _ {2} ^ {*}.
\]

The real number \((\pmb{u}|\pmb{x} + \pmb{y})\) is an integer for all \(\pmb{x} \in \mathbf{G}_1\) and all \(\pmb{y} \in \mathbf{G}_2\) if and only if \((\pmb{u}|\pmb{x})\) is an integer for all \(\pmb{x} \in \mathbf{G}_1\) and \((\pmb{u}|\pmb{y})\) is an integer for all \(\pmb{y} \in \mathbf{G}_2\) , because \((\pmb{u}|\pmb{x} + \pmb{y}) = (\pmb{u}|\pmb{x}) + (\pmb{u}|\pmb{y})\) ; hence

\[
(\mathrm{G} _ {1} + \mathrm{G} _ {2}) ^ {*} = \mathrm{G} _ {1} ^ {*} \cap \mathrm{G} _ {2} ^ {*}
\]

for any two subgroups \(G_{1}\) , \(G_{2}\) of \(R^{n}\) . If now we suppose that \(G_{1}\) and \(G_{2}\) are closed, we have \((\mathbf{G}_{1}^{*} + \mathbf{G}_{2}^{*})^{*} = \mathbf{G}_{1} \cap \mathbf{G}_{2}\) by Proposition 6; hence, taking the associated subgroups and applying Proposition 6 again, \((\mathbf{G}_{1} \cap \mathbf{G}_{2})^{*} = \overline{\mathbf{G}_{1}^{*} + \mathbf{G}_{2}^{*}}\) .

Remark. Let \(\mathbf{G}_1, \mathbf{G}_2\) be two lattices in \(\mathbf{R}^n\) (no. 1) such that \(\mathbf{G}_2 \subset \mathbf{G}_1\) ; then (Proposition 5) \(\mathbf{G}_1^*\) and \(\mathbf{G}_2^*\) are lattices in \(\mathbf{R}^n\) such that \(\mathbf{G}_1^* \subset \mathbf{G}_2^*\) . We have seen in no. 1 that there is an integer \(m > 0\) such that \(m\mathbf{G}_1 \subset \mathbf{G}_2\) ;

consequently for \(x \in G_1\) and \(u \in G_2^*\) we have \(m(u|x) \in Z\) and therefore \((u|x) \in Q\) . If \(x \in G_2\) and \(u \in G_2^*\) , or if \(x \in G_1\) and \(u \in G_1^*\) , we have \((u|x) \in Z\) by definition. Consequently on passing to the quotients the \(Z\) -bilinear mapping \((x, u) \to (u|x)\) of \(G_1 \times G_2^*\) into \(Q\) defines a \(Z\) -bilinear mapping \(B\) of \((G_1/G_2) \times (G_2^*/G_1^*)\) into \(Q/Z\) . Moreover, it is clear that if \(\bar{x}_0 \in G_1/G_2\) (resp. \(\bar{u}_0 \in G_2^*/G_1^*\) ) is such that, for each \(\bar{u} \in G_2^*/G_1^*\) (resp. for each \(\bar{x} \in G_1/G_2\) ) we have

\[
\mathrm{B} (\overline {{{\boldsymbol {x}}}} _ {0}, \boldsymbol {u}) = \mathrm{o} \quad [ \text { resp. } \mathrm{B} (\overline {{{\boldsymbol {x}}}}, \overline {{{\boldsymbol {u}}}} _ {0}) = \mathrm{o} ],
\]

then necessarily \(\overline{\mathbf{x}}_0 = \mathbf{o}\) (resp. \(\overline{\mathbf{u}}_0 = \mathbf{o}\) ). It follows that there is a \(\mathbf{Z}\) -linear bijection \(h\) of \(\mathrm{G}_2^* / \mathrm{G}_1^*\) onto the dual of \(\mathrm{G}_1 / \mathrm{G}_2\) , such that \(\langle \overline{\mathbf{x}}, h(\overline{\mathbf{u}}) \rangle = \mathrm{B}(\overline{\mathbf{x}}, \overline{\mathbf{u}})\) for \(\overline{\mathbf{x}} \in \mathrm{G}_1 / \mathrm{G}_2\) and \(\overline{\mathbf{u}} \in \mathrm{G}_2^* / \mathrm{G}_1^*\) ; in particular, the finite groups \(\mathrm{G}_1 / \mathrm{G}_2\) and \(\mathrm{G}_2^* / \mathrm{G}_1^*\) are isomorphic.

